\section{Dados coletados}
Esta seção visa apenas apresentar e definir quais dados foram coletados. A próxima seção descreve com mais detalhes como a coleta foi realizada. Para cada combinação de algoritmo, tamanho e tipo de vetor, foi registrado o número de comparações, o número de movimentações e o tempo de execução. Essas três métricas são definidas na tabela \ref{tab:dados-coletados}.

Obter essas informações de cada algoritmo é importante para sabermos na prática como eles se comportam para diferentes configurações de entrada. Por exemplo, um algoritmo que executa muitas movimentações pode não ser adequado em um cenário em que mudar os registros de posição seja uma operação muito lenta.

\begin{table}[H]
    \centering
    \Caption{\label{tab:dados-coletados}Dados coletados e suas definições}
    \begin{tabular}{ c | l }
        \hline
        Dado                    & Definição\\
        \hline
        Número de comparações   & \makecell[l]{Uma comparação ocorre quando um elemento do\\
                                               vetor é comparado com outro elemento do vetor ou\\
                                               com outra variável do programa}\\
        \hline
        Número de movimentações & \makecell[l]{Uma movimentação ocorre sempre que um valor é\\
                                               atribuído a um elemento do vetor ou um elemento\\
                                               do vetor é atribuído a uma outra variável}\\
        \hline
        Tempo de execução       & \makecell[l]{Tempo decorrido entre o início e o fim da\\
                                               execução do algoritmo em milissegundos}\\
        \hline
    \end{tabular}
\end{table}
